\section{Conclusion}\label{sec:conclusion}

Networked XR has grown almost fourfold in eight years and has changed character on the way. It began as
a multimedia-systems field centred on 360-degree streaming. Its claims were tested with real players,
real traces and, often, real viewers. Today it is dominated by resource allocation, edge computing and 6G
architecture. It models XR traffic generically and tests most new claims in simulation. The empirical
core has not disappeared, but it has not grown with the rest. This core covers latency measurement,
multi-user and holographic systems, and QoE studies.

Organising the field by function makes the consequences visible. The functions that build systems,
compute and connect, rest on real testbeds in \fnComputeReal{} and \fnConnectReal{} of their papers.
Provisioning, now the largest and fastest-growing function, does so in \fnProvisionReal{}. Across all six
functions the same pattern recurs. Optimisers are well posed, but their inputs are assumed rather than
measured. The QoE they maximise was never shown to users. The requirements they meet were copied from
earlier surveys. The traffic they carry was synthesised.

The remedy is less a new algorithm than shared ground: open, frame-synchronised traces of real XR
workloads, and a QoE model fitted to people experiencing replayed network dynamics. With these two
artefacts, the questions that recur in every section become answerable on common terms: whether a learned
scheduler beats a heuristic on a real cell, how much latency a defence costs, and what a hologram really
needs. Without them, the field will keep growing faster than its evidence.
